\subsubsection{Word Similarity Ignorance}\label{sec:word-similarity-ignorance}

The previous section showed that traditional representations such as one-hot encoding are sparse and high-dimensional. However, an even more serious problem exists: \textbf{one-hot vectors completely ignore Semantic Similarity between words}.

\highspace
\begin{definitionbox}[: Semantic Similarity]\label{def:semantic-similarity}
    \definition{Semantic Similarity} is the \hl{degree to which two words, phrases, or concepts have similar meanings.} More formally, two linguistic entities are semantically similar if they can be used in similar contexts and convey related meanings. 

    \highspace
    In the context of Word Embeddings, words with similar meanings should have similar vectors, then:
    \begin{equation*}
        \text{semantic similarity} \, \Rightarrow \, \text{geometric proximity}
    \end{equation*}
    The whole goal of word embeddings is to transform an abstract concept such as \textbf{meaning} into something measurable using geometry.

    \highspace
    \textcolor{Red2}{\faIcon{exclamation-triangle} \textbf{Similarity vs Relatedness.}} There is an important nuance. Consider \emph{car} and \emph{road}. They are strongly related, but they do not mean the same thing. Similarly, \emph{doctor} and \emph{hospital} are related, but not similar. In contrast, \emph{car} and \emph{automobile} are both related and similar. Therefore, many embedding models actually capture a mixture of:
    \begin{itemize}
        \item \textbf{Semantic Similarity} is the degree to which two words have similar meanings and can be used interchangeably in many contexts. For example, \emph{car} and \emph{automobile} are semantically similar because they refer to the same concept.
        \item \definition{Semantic Relatedness} is the degree to which two words are associated or connected, regardless of whether they have similar meanings. For example, \emph{car} and \emph{road} are semantically related because they frequently occur together and belong to the same conceptual domain, but they are not semantically similar because they refer to different concepts.
    \end{itemize}
\end{definitionbox}

\highspace
For example, consider the following twosentences:
\begin{itemize}
    \item \emph{Obama speaks to the media in Illinois.}
    \item \emph{The President addresses the press in Chicago.}
\end{itemize}
To a human reader, these sentences convey almost the same meaning. The words are different, but the concepts they express are strongly related:
\begin{center}
    \begin{tabular}{@{} l l @{}}
        \toprule
        \textbf{Sentence 1} & \textbf{Sentence 2} \\
        \midrule
        Obama & President \\
        speaks & addresses \\
        media & press \\
        Illinois & Chicago \\
        \bottomrule
    \end{tabular}
\end{center}

\noindent
A human immediately recognizes these semantic relationships. A machine using one-hot representations does not because the one-hot vectors for these words are orthogonal to each other. In other words, the one-hot representation does not capture any notion of similarity between words. Let's see this in action.

\highspace
\begin{flushleft}
    \textcolor{Red2}{\faIcon{exclamation-triangle} \textbf{One-Hot Encoding Cannot Express Similarity}}
\end{flushleft}
Recall that in a vocabulary of size $V$, each word is represented by a vector containing a single $1$ and $V-1$ zeros. For example:
\begin{align*}
    \text{speaks} &= [0, 0, 1, 0, \dots, 0] \\
    \text{addresses} &= [0, 0, 0, 0, \dots, 1, 0]
\end{align*}
Similarly:
\begin{align*}
    \text{obama} &= [0, 0, 0, 0, \dots, 1, 0, 0] \\
    \text{president} &= [0, 0, 0, 1, \dots, 0, 0, 0]
\end{align*}
And:
\begin{align*}
    \text{illinois} &= [1, 0, 0, 0, \dots, 0] \\
    \text{chicago} &= [0, 1, 0, 0, \dots, 0]
\end{align*}
These vectors have no overlapping active components (i.e., no $1$s in the same position), so they are orthogonal.

\highspace
Mathematically, \textbf{any two different one-hot vectors are orthogonal}. For example:
\begin{align*}
    \text{speaks}^{T} \cdot \text{addresses} &= 0 \\
    \text{obama}^{T} \cdot \text{president} &= 0 \\
    \text{illinois}^{T} \cdot \text{chicago} &= 0
\end{align*}
We can also use the symbol $\perp$ to denote orthogonality:
\begin{align*}
    \text{speaks} \perp \text{addresses} \\
    \text{obama} \perp \text{president} \\
    \text{illinois} \perp \text{chicago}
\end{align*}
Geometrically, the angle between the vectors is $90^{\circ}$, which means they are completely unrelated.

\highspace
\textcolor{Red2}{\faIcon{question-circle} \textbf{Why is this a problem?}} Suppose a language model has learned many examples containing the word ``\emph{president}'' and later encounters ``\emph{Obama}''. A human immediately understands that these words are related, but the \textbf{model cannot}. From the model's perspective ``\emph{Obama}'' is just another compleltely unrelated symbol. And the same happens for ``\emph{speaks}'' and ``\emph{addresses}'', or ``\emph{Illinois}'' and ``\emph{Chicago}''. Since the \hl{vectors are orthogonal}, the \hl{model has no mechanism for transferring knowledge from one word to another}.

\highspace
\begin{flushleft}
    \textcolor{Red2}{\faIcon{times-circle} \textbf{Failure to Generalize}}
\end{flushleft}
This inability to capture similarity directly \textbf{affects a model's capacity to generalize}. Imagine that during training we only observe: ``\emph{Obama speaks to the media in Illinois}''. If we later encounter: ``\emph{The President addresses the press in Chicago}'' the model should ideally recognize that the second sentence is semantically similar to the first.

\highspace
With one-hot representations, however, every key word is represented by a completely different orthogonal vector. Consequently, the model sees the second sentence as almost entirely new. The knowledge learned from the first sentence cannot be reused effectively.

\highspace
\begin{flushleft}
    \textcolor{Green3}{\faIcon{check-circle} \textbf{What we actually need: Word Embeddings}}
\end{flushleft}
To generalize successfully, \textbf{similar words} should have \textbf{similar representations}. Ideally, we would like:
\begin{equation*}
    \text{similar meaning} \, \Rightarrow \, \text{similar vectors}
\end{equation*}
Word pairs share no similarity in their one-hot representations, but we need word similarity to be reflected in the representations.

\highspace
The fundamental \textbf{goal} of \textbf{word embeddings} is therefore to \hl{replace orthogonal one-hot vectors with dense continuous vectors} such that:
\begin{align*}
    \mathbf{v}_{\text{Obama}} &\approx \mathbf{v}_{\text{President}} \\
    \mathbf{v}_{\text{speaks}} &\approx \mathbf{v}_{\text{addresses}} \\
    \mathbf{v}_{\text{Illinois}} &\approx \mathbf{v}_{\text{Chicago}}
\end{align*}
Instead of \textbf{representing words} as isolated symbols, we want to represent them \textbf{as points in a continuous semantic space where distance reflects meaning}. This idea leads directly to the concept of word embeddings, which we introduce next.